% ============================================================
% Question Bank: ch02-09 Financial Literacy --- Budgeting, Saving, and Investing
% Topic Title: Financial Literacy --- Budgeting, Saving, and Investing
% CCSS: Supplementary
% Questions: 30 (18 MC + 7 SA + 5 Graphical)
% ============================================================

% --- Multiple Choice Questions (18) ---

\begin{practiceQuestion}{2-9-q01}{mc}
\begin{questionText}
A budget shows income $= \$1{,}800$, fixed expenses $= \$1{,}000$, variable expenses $= \$500$. How much is left to save?
\end{questionText}
\choiceA{$\$200$}
\choiceB{$\$300$}
\choiceC{$\$500$}
\choiceD{$\$800$}
\correctAnswer{B}
\explanation{Savings $= 1{,}800 - 1{,}000 - 500 = \$300$.}
\end{practiceQuestion}

\begin{practiceQuestion}{2-9-q02}{mc}
\begin{questionText}
Which of the following is a fixed expense?
\end{questionText}
\choiceA{Groceries}
\choiceB{Movie tickets}
\choiceC{Monthly rent}
\choiceD{Clothing}
\correctAnswer{C}
\explanation{Rent stays the same each month, making it a fixed expense. The others vary.}
\end{practiceQuestion}

\begin{practiceQuestion}{2-9-q03}{mc}
\begin{questionText}
You earn $\$2{,}400$ per month and save $\$240$. What percent of your income do you save?
\end{questionText}
\choiceA{$24\%$}
\choiceB{$12\%$}
\choiceC{$10\%$}
\choiceD{$20\%$}
\correctAnswer{C}
\explanation{$\frac{240}{2{,}400} = 0.10 = 10\%$.}
\end{practiceQuestion}

\begin{practiceQuestion}{2-9-q04}{mc}
\begin{questionText}
Which is an example of a variable expense?
\end{questionText}
\choiceA{Car payment}
\choiceB{Phone plan}
\choiceC{Rent}
\choiceD{Entertainment}
\correctAnswer{D}
\explanation{Entertainment costs change from month to month, making it a variable expense.}
\end{practiceQuestion}

\begin{practiceQuestion}{2-9-q05}{mc}
\begin{questionText}
You want to save $\$900$ in $6$ months. How much must you save each month?
\end{questionText}
\choiceA{$\$100$}
\choiceB{$\$150$}
\choiceC{$\$180$}
\choiceD{$\$200$}
\correctAnswer{B}
\explanation{$900 \div 6 = \$150$ per month.}
\end{practiceQuestion}

\begin{practiceQuestion}{2-9-q06}{mc}
\begin{questionText}
Investing means putting money into stocks or bonds. Compared to a savings account, investing typically has:
\end{questionText}
\choiceA{Lower risk and lower potential growth}
\choiceB{Lower risk and higher potential growth}
\choiceC{Higher risk and higher potential growth}
\choiceD{No risk and guaranteed growth}
\correctAnswer{C}
\explanation{Investments can earn more than a savings account, but you could also lose money.}
\end{practiceQuestion}

\begin{practiceQuestion}{2-9-q07}{mc}
\begin{questionText}
Your monthly income is $\$1{,}600$. Fixed expenses are $\$900$ and variable expenses are $\$450$. What is your savings rate?
\end{questionText}
\choiceA{$6.25\%$}
\choiceB{$15.625\%$}
\choiceC{$10\%$}
\choiceD{$25\%$}
\correctAnswer{B}
\explanation{Savings $= 1{,}600 - 900 - 450 = \$250$. Rate: $\frac{250}{1{,}600} = 0.15625 = 15.625\%$.}
\end{practiceQuestion}

\begin{practiceQuestion}{2-9-q08}{mc}
\begin{questionText}
A savings account earns $2\%$ per year. You deposit $\$500$ for $1$ year. How much interest do you earn?
\end{questionText}
\choiceA{$\$2$}
\choiceB{$\$10$}
\choiceC{$\$50$}
\choiceD{$\$100$}
\correctAnswer{B}
\explanation{Interest $= 500 \times 0.02 = \$10$.}
\end{practiceQuestion}

\begin{practiceQuestion}{2-9-q09}{mc}
\begin{questionText}
Your food budget is $\$300$ per month. You spend $\$330$ one month. By what percent did you go over budget?
\end{questionText}
\choiceA{$10\%$}
\choiceB{$3\%$}
\choiceC{$30\%$}
\choiceD{$11\%$}
\correctAnswer{A}
\explanation{Over by $\$30$. Percent: $\frac{30}{300} = 0.10 = 10\%$.}
\end{practiceQuestion}

\begin{practiceQuestion}{2-9-q10}{mc}
\begin{questionText}
Experts suggest saving at least what percent of your income?
\end{questionText}
\choiceA{$1\%$}
\choiceB{$5\%$}
\choiceC{$10\%$}
\choiceD{$50\%$}
\correctAnswer{C}
\explanation{Financial experts generally recommend saving at least $10\%$ of income.}
\end{practiceQuestion}

\begin{practiceQuestion}{2-9-q11}{mc}
\begin{questionText}
Income: $\$3{,}000$. After all expenses, you save $\$450$. What are your total expenses?
\end{questionText}
\choiceA{$\$3{,}450$}
\choiceB{$\$2{,}550$}
\choiceC{$\$2{,}450$}
\choiceD{$\$3{,}000$}
\correctAnswer{B}
\explanation{Expenses $= 3{,}000 - 450 = \$2{,}550$.}
\end{practiceQuestion}

\begin{practiceQuestion}{2-9-q12}{mc}
\begin{questionText}
Which is the BEST reason to choose a savings account over investing?
\end{questionText}
\choiceA{It always earns more money}
\choiceB{Your money is safe and protected}
\choiceC{You never pay fees}
\choiceD{You can withdraw unlimited times}
\correctAnswer{B}
\explanation{Savings accounts are low-risk — your money is safe, unlike investments that could lose value.}
\end{practiceQuestion}

\begin{practiceQuestion}{2-9-q13}{mc}
\begin{questionText}
Your entertainment budget is $\$80$/month. After $5$ months, you've spent $\$350$ total on entertainment. Are you under or over budget?
\end{questionText}
\choiceA{Under budget by $\$50$}
\choiceB{Over budget by $\$50$}
\choiceC{Under budget by $\$30$}
\choiceD{Over budget by $\$30$}
\correctAnswer{A}
\explanation{Budget for $5$ months: $80 \times 5 = \$400$. Spent: $\$350$. Under by $400 - 350 = \$50$.}
\end{practiceQuestion}

\begin{practiceQuestion}{2-9-q14}{mc}
\begin{questionText}
You deposit $\$1{,}200$ in a savings account at $3\%$ annual interest. After $1$ year, your balance is:
\end{questionText}
\choiceA{$\$1{,}203$}
\choiceB{$\$1{,}236$}
\choiceC{$\$1{,}260$}
\choiceD{$\$1{,}236$}
\correctAnswer{B}
\explanation{Interest $= 1{,}200 \times 0.03 = \$36$. Balance $= 1{,}200 + 36 = \$1{,}236$.}
\end{practiceQuestion}

\begin{practiceQuestion}{2-9-q15}{mc}
\begin{questionText}
A student earns $\$200$ per month from a part-time job. She spends $\$60$ on a phone plan (fixed) and $\$80$ on food (variable). How much can she save?
\end{questionText}
\choiceA{$\$140$}
\choiceB{$\$60$}
\choiceC{$\$80$}
\choiceD{$\$120$}
\correctAnswer{B}
\explanation{$200 - 60 - 80 = \$60$ available to save.}
\end{practiceQuestion}

\begin{practiceQuestion}{2-9-q16}{mc}
\begin{questionText}
An investment averages $8\%$ return per year. A savings account earns $2\%$ per year. On a $\$1{,}000$ deposit for $1$ year, how much more does the investment earn (if it meets its average)?
\end{questionText}
\choiceA{$\$6$}
\choiceB{$\$60$}
\choiceC{$\$80$}
\choiceD{$\$100$}
\correctAnswer{B}
\explanation{Investment: $1{,}000 \times 0.08 = \$80$. Savings: $1{,}000 \times 0.02 = \$20$. Difference: $80 - 20 = \$60$.}
\end{practiceQuestion}

\begin{practiceQuestion}{2-9-q17}{mc}
\begin{questionText}
You budget $\$500$ for groceries each month. In March, you spend $\$475$. In April, you spend $\$540$. What is the average monthly spending?
\end{questionText}
\choiceA{$\$500$}
\choiceB{$\$507.50$}
\choiceC{$\$515$}
\choiceD{$\$540$}
\correctAnswer{B}
\explanation{Average: $(475 + 540) \div 2 = 1{,}015 \div 2 = \$507.50$.}
\end{practiceQuestion}

\begin{practiceQuestion}{2-9-q18}{mc}
\begin{questionText}
Which formula correctly represents the relationship between income, expenses, and savings?
\end{questionText}
\choiceA{Income $=$ Savings $+$ Expenses}
\choiceB{Savings $=$ Expenses $-$ Income}
\choiceC{Income $=$ Expenses $-$ Savings}
\choiceD{Expenses $=$ Income $+$ Savings}
\correctAnswer{A}
\explanation{Income $-$ Expenses $=$ Savings, which can be rearranged to Income $=$ Savings $+$ Expenses.}
\end{practiceQuestion}

% --- Short Answer Questions (7) ---

\begin{practiceQuestion}{2-9-q19}{sa}
\begin{questionText}
Monthly income: $\$2{,}500$. Fixed expenses: $\$1{,}400$. Variable expenses: $\$700$. How much is left to save?
\end{questionText}
\correctAnswer{$\$400$}
\explanation{$2{,}500 - 1{,}400 - 700 = \$400$.}
\end{practiceQuestion}

\begin{practiceQuestion}{2-9-q20}{sa}
\begin{questionText}
You want to save $\$1{,}800$ in $12$ months. How much must you save per month?
\end{questionText}
\correctAnswer{$\$150$}
\explanation{$1{,}800 \div 12 = \$150$ per month.}
\end{practiceQuestion}

\begin{practiceQuestion}{2-9-q21}{sa}
\begin{questionText}
Your clothing budget is $\$200$ per month. You spend $\$230$. By what percent did you go over budget?
\end{questionText}
\correctAnswer{$15\%$}
\explanation{Over by $\$30$. Percent: $\frac{30}{200} = 0.15 = 15\%$.}
\end{practiceQuestion}

\begin{practiceQuestion}{2-9-q22}{sa}
\begin{questionText}
You earn $\$1{,}000$ per month. You save $\$125$. What is your savings rate as a percent?
\end{questionText}
\correctAnswer{$12.5\%$}
\explanation{$\frac{125}{1{,}000} = 0.125 = 12.5\%$.}
\end{practiceQuestion}

\begin{practiceQuestion}{2-9-q23}{sa}
\begin{questionText}
A savings account earns $4\%$ per year. You deposit $\$750$ for $1$ year. How much interest do you earn?
\end{questionText}
\correctAnswer{$\$30$}
\explanation{$I = 750 \times 0.04 = \$30$.}
\end{practiceQuestion}

\begin{practiceQuestion}{2-9-q24}{sa}
\begin{questionText}
Name one advantage of investing and one advantage of a savings account.
\end{questionText}
\correctAnswer{Investing: higher potential growth; Savings: low risk}
\explanation{Investments can grow faster but carry risk. Savings accounts are safe but earn less.}
\end{practiceQuestion}

\begin{practiceQuestion}{2-9-q25}{sa}
\begin{questionText}
You earn $\$3{,}200$ per month. Your goal is to save $15\%$ of income. How many dollars should you save each month?
\end{questionText}
\correctAnswer{$\$480$}
\explanation{$3{,}200 \times 0.15 = \$480$.}
\end{practiceQuestion}

% --- Graphical Questions (3 GMC + 2 GSA) ---

\begin{practiceQuestion}{2-9-q26}{gmc}
\begin{questionText}
The circle graph shows how a student spends her monthly income of $\$800$. How much does she spend on food?

\begin{center}
\begin{tikzpicture}[scale=1.2]
  % Pie chart
  \fill[funBlue!60] (0,0) -- (90:1.5) arc (90:198:1.5) -- cycle;
  \fill[funRed!60] (0,0) -- (198:1.5) arc (198:306:1.5) -- cycle;
  \fill[green!50!black] (0,0) -- (306:1.5) arc (306:378:1.5) -- cycle;
  \fill[orange!60] (0,0) -- (18:1.5) arc (18:90:1.5) -- cycle;
  \draw[thick] (0,0) circle (1.5);
  % Labels
  \node at (144:1) {\footnotesize Rent};
  \node at (144:0.55) {\footnotesize $30\%$};
  \node at (252:1) {\footnotesize Food};
  \node at (252:0.55) {\footnotesize $30\%$};
  \node at (342:1) {\footnotesize Fun};
  \node at (342:0.55) {\footnotesize $20\%$};
  \node at (54:1) {\footnotesize Save};
  \node at (54:0.55) {\footnotesize $20\%$};
\end{tikzpicture}
\end{center}
\end{questionText}
\choiceA{$\$160$}
\choiceB{$\$200$}
\choiceC{$\$240$}
\choiceD{$\$300$}
\correctAnswer{C}
\explanation{Food is $30\%$ of $\$800$: $800 \times 0.30 = \$240$.}
\end{practiceQuestion}

\begin{practiceQuestion}{2-9-q27}{gmc}
\begin{questionText}
The bar graph shows a family's monthly budget. Which category is a variable expense?

\begin{center}
\begin{tikzpicture}[scale=0.6]
  \draw[thick,->] (0,0) -- (0,6) node[above] {\$};
  \draw[thick,->] (0,0) -- (10,0);
  \foreach \y/\label in {0/0, 1/200, 2/400, 3/600, 4/800, 5/1000} {
    \draw (-0.1,\y) -- (0.1,\y);
    \node[left] at (-0.15,\y) {\footnotesize \$\label};
  }
  \fill[funBlue!70] (0.5,0) rectangle (2,5);
  \node[below] at (1.25,0) {\footnotesize Rent};
  \node[above] at (1.25,5) {\footnotesize \$1000};
  \fill[funRed!70] (3,0) rectangle (4.5,2);
  \node[below] at (3.75,0) {\footnotesize Groceries};
  \node[above] at (3.75,2) {\footnotesize \$400};
  \fill[green!60] (5.5,0) rectangle (7,1.5);
  \node[below] at (6.25,0) {\footnotesize Phone};
  \node[above] at (6.25,1.5) {\footnotesize \$300};
  \fill[orange!70] (8,0) rectangle (9.5,1);
  \node[below] at (8.75,0) {\footnotesize Fun};
  \node[above] at (8.75,1) {\footnotesize \$200};
\end{tikzpicture}
\end{center}
\end{questionText}
\choiceA{Rent}
\choiceB{Phone}
\choiceC{Groceries}
\choiceD{All of them}
\correctAnswer{C}
\explanation{Groceries is a variable expense because the amount you spend on food changes month to month. Rent and phone plans are fixed amounts.}
\end{practiceQuestion}

\begin{practiceQuestion}{2-9-q28}{gmc}
\begin{questionText}
The line graph shows a student's savings account balance over $4$ months. Between which two months did the balance increase the most?

\begin{center}
\begin{tikzpicture}[scale=0.8]
  \draw[thick,->] (0,0) -- (5.5,0) node[right] {Month};
  \draw[thick,->] (0,0) -- (0,5) node[above] {\$};
  \foreach \x/\label in {1/Jan, 2/Feb, 3/Mar, 4/Apr} {
    \draw (\x,-0.1) -- (\x,0.1);
    \node[below] at (\x,-0.15) {\footnotesize \label};
  }
  \foreach \y/\label in {1/100, 2/200, 3/300, 4/400} {
    \draw (-0.1,\y) -- (0.1,\y);
    \node[left] at (-0.15,\y) {\footnotesize \$\label};
  }
  \draw[thick,funBlue,mark=*] plot coordinates {(1,1) (2,1.5) (3,3) (4,3.5)};
  \node[above] at (1,1) {\footnotesize \$100};
  \node[above] at (2,1.5) {\footnotesize \$150};
  \node[above] at (3,3) {\footnotesize \$300};
  \node[above right] at (4,3.5) {\footnotesize \$350};
\end{tikzpicture}
\end{center}
\end{questionText}
\choiceA{Jan to Feb}
\choiceB{Feb to Mar}
\choiceC{Mar to Apr}
\choiceD{Jan to Mar}
\correctAnswer{B}
\explanation{Jan to Feb: $\$50$ increase. Feb to Mar: $\$150$ increase. Mar to Apr: $\$50$ increase. The biggest increase is Feb to Mar.}
\end{practiceQuestion}

\begin{practiceQuestion}{2-9-q29}{gsa}
\begin{questionText}
The table shows a monthly budget. Find the savings rate (as a percent of income).

\begin{center}
\begin{tabular}{|l|c|}
\hline
\textbf{Category} & \textbf{Amount} \\
\hline
Income & $\$2{,}000$ \\
\hline
Rent (fixed) & $\$700$ \\
\hline
Phone (fixed) & $\$100$ \\
\hline
Groceries (variable) & $\$350$ \\
\hline
Entertainment (variable) & $\$200$ \\
\hline
Transportation (variable) & $\$150$ \\
\hline
\end{tabular}
\end{center}
\end{questionText}
\correctAnswer{$25\%$}
\explanation{Total expenses: $700 + 100 + 350 + 200 + 150 = \$1{,}500$. Savings: $2{,}000 - 1{,}500 = \$500$. Rate: $\frac{500}{2{,}000} = 0.25 = 25\%$.}
\end{practiceQuestion}

\begin{practiceQuestion}{2-9-q30}{gsa}
\begin{questionText}
The bar graph shows two people's monthly budgets. Who saves more money, and by how much?

\begin{center}
\begin{tikzpicture}[scale=0.55]
  \draw[thick,->] (0,0) -- (0,7) node[above] {\$};
  \draw[thick,->] (0,0) -- (12,0);
  \foreach \y/\label in {0/0, 1/200, 2/400, 3/600, 4/800, 5/1000, 6/1200} {
    \draw (-0.1,\y) -- (0.1,\y);
    \node[left] at (-0.15,\y) {\footnotesize \$\label};
  }
  % Alex
  \fill[funBlue!50] (0.5,0) rectangle (1.8,5);
  \node[below] at (1.15,0) {\footnotesize Income};
  \node[above] at (1.15,5) {\footnotesize \$1000};
  \fill[funRed!50] (2.2,0) rectangle (3.5,4);
  \node[below] at (2.85,0) {\footnotesize Expenses};
  \node[above] at (2.85,4) {\footnotesize \$800};
  \node[below] at (2,-0.8) {\footnotesize Alex};

  % Beth
  \fill[funBlue!50] (6,0) rectangle (7.3,6);
  \node[below] at (6.65,0) {\footnotesize Income};
  \node[above] at (6.65,6) {\footnotesize \$1200};
  \fill[funRed!50] (7.7,0) rectangle (9,5.25);
  \node[below] at (8.35,0) {\footnotesize Expenses};
  \node[above] at (8.35,5.25) {\footnotesize \$1050};
  \node[below] at (7.5,-0.8) {\footnotesize Beth};
\end{tikzpicture}
\end{center}
\end{questionText}
\correctAnswer{Alex by $\$50$}
\explanation{Alex saves: $1{,}000 - 800 = \$200$. Beth saves: $1{,}200 - 1{,}050 = \$150$. Alex saves $200 - 150 = \$50$ more.}
\end{practiceQuestion}
